% Generated draft: requires mathematical and source review.
\documentclass[11pt]{article}
\usepackage[margin=25mm]{geometry}
\usepackage{fontspec}
\setmainfont{Latin Modern Roman}
\usepackage{amsmath,amssymb,mathtools,unicode-math}
\setmathfont{Latin Modern Math}
\usepackage{enumitem,needspace,xurl,hyperref}
\setlength{\parindent}{0pt}
\setlength{\parskip}{5pt}
\setlength{\emergencystretch}{3em}
\newcommand{\N}{\mathbb{N}}
\newcommand{\Z}{\mathbb{Z}}
\newcommand{\Q}{\mathbb{Q}}
\newcommand{\R}{\mathbb{R}}
\newcommand{\C}{\mathbb{C}}
\newcommand{\F}{\mathbb{F}}
\newcommand{\A}{\mathbb{A}}
\newcommand{\PP}{\mathbb{P}}
\newcommand{\E}{\mathbb{E}}
\newcommand{\eps}{\varepsilon}
\newcommand{\dd}{\,\mathrm{d}}
\begin{document}
\section*{UnsolvedMath ID: 30002288}
\textbf{Source problem code:} catalogue entry
\subsection{Definitions and mathematical statement}
The supplied record poses two questions about first fractional infinity eigenfunctions. It does not specify the underlying domain, the admissible exponent range, or the eigenvalue equations. Consequently, the following is a conditional statement of the recorded questions, requiring the missing specifications before it defines a complete mathematical problem.

Let \(\Omega\) denote the intended domain, let \(\alpha\) denote the intended exponent, and let \(\delta\) and \(\Gamma\) denote the geometric data appearing in the source. The supplied verification note calls \(\Gamma\) a high-ridge set, but the record does not reproduce its definition or the definition of \(\delta\). For a nonempty subset \(A\) of the intended ambient metric space, write
\[
\operatorname{dist}(x,A)=\inf_{a\in A}d(x,a),
\]
where \(d\) is the intended ambient distance, which is not specified in the record. For \(\Gamma_1\subseteq\Gamma\), the proposed representation is
\[
\varrho_1(x)=\operatorname{dist}(x,\Gamma_1),\qquad
u_{\Gamma_1}(x)=\frac{\delta(x)^\alpha}{\delta(x)^\alpha+\varrho_1(x)^\alpha}.
\]
This expression applies only where its real powers are defined and its denominator is nonzero. The record does not specify admissibility conditions on \(\Gamma_1\), including whether it must be nonempty or closed, or how boundary values are assigned.

\emph{Representation question.} With the domain, operator, boundary condition, solution concept, and normalization fixed as intended by the cited source, is it true that for every first fractional infinity eigenfunction \(u\), there exists an admissible subset \(\Gamma_1\subseteq\Gamma\) such that
\[
u(x)=u_{\Gamma_1}(x)\qquad(x\in\Omega)?
\]
The normalization qualification is unresolved: the record asks for every first eigenfunction but does not specify how scalar multiples are treated.

\emph{Limit question.} Set
\[
u_{\Gamma}(x)=\frac{\delta(x)^\alpha}{\delta(x)^\alpha+\operatorname{dist}(x,\Gamma)^\alpha}.
\]
The source calls this the maximal solution. Is \(u_{\Gamma}\) the limit, as \(p\to\infty\), of the intended fractional \(p\)-eigenfunctions? The record does not specify their equations, selection, normalization, convergence topology, or whether convergence concerns every choice, a particular family, or a subsequence. No choice among these formulations is made here. The term maximal is retained as supplied terminology; its comparison class and ordering require confirmation.
\subsection{Short English statement}
Does the displayed distance formula, with a suitable subset of the high ridge, represent every appropriately normalized first fractional infinity eigenfunction? Do fractional \(p\)-eigenfunctions converge as \(p\to\infty\) to the solution obtained from the whole high ridge? The supplied record leaves essential geometric, analytic, normalization, and convergence specifications unstated.
\subsection{Sources}
\begin{itemize}
\item Supplied local UnsolvedMath registry record: permanent ID \(30002288\), exact source code \texttt{catalogue entry}, title \emph{Representation and Limits of Fractional Infinity Eigenfunctions}, Oberwolfach Reports Open Problems collection. The record gives proposed year 2013 and names no individual proposer.
\item \emph{Mini-Workshop: The p-Laplacian Operator and Applications} (2013), Oberwolfach Reports, pp.\ 433--482, DOI \texttt{10.4171/owr/2013/08}. Supplied report reference and literature-source link: \url{https://doi.org/10.4171/owr/2013/08}.
\item Lindgren--Lindqvist, \emph{Fractional Eigenvalues}. Supplied literature-source link: \url{https://doi.org/10.1007/s00526-013-0600-1}.
\item Literature assessment dated exactly \texttt{2026-08-21}: the supplied curation labels the problem open and states that Lindgren--Lindqvist explicitly leave open both completeness of the \(\Gamma_1\)-representation and convergence of fractional \(p\)-eigenfunctions to the maximal solution. This is supplied dated curation derived from the Oberwolfach material, not an independently verified literature assessment; time-sensitive status and cited claims require independent verification.
\item The registry labels the statement \texttt{corrected\_verified} and reports that the official report explicitly poses both questions and supplies the high-ridge-subset and distance-function definitions. Those definitions are not reproduced in the supplied record. This verification claim is preserved as registry metadata, not adopted as independent mathematical verification.
\end{itemize}
\end{document}
